\section{Introduction}

The exponential proliferation of high-speed digital networks and over-the-top (OTT) media services has fundamentally transformed global multimedia distribution. High-definition video content represents highly valuable intellectual property for creators and distribution networks alike. Consequently, digital assets have become primary targets for widespread online piracy, unauthorized redistribution, and man-in-the-middle interception \cite{bhattacharyya2021multimedia}. To combat these copyright violations, content providers rely heavily on Digital Rights Management (DRM) frameworks to enforce robust access control, manage licensing compliance, and maintain end-to-end data confidentiality.

Standard commercial DRM ecosystems, such as Google Widevine, Apple FairPlay, and Microsoft PlayReady, offer industrial-grade security but suffer from severe architectural rigidity \cite{nussbaumer2022drm}. These industrial platforms typically rely on full-stream cryptographic encapsulation, where every single byte of the compressed video bitstream payload must pass through symmetric decryption pipelines at the client-side media rendering engine. While secure, executing full-payload decryption inflicts intensive computational overhead, elevated battery depletion rates, and micro-stuttering latency spikes on resource-constrained endpoint devices \cite{subramanyam2023lightweight}. These constraints heavily impact low-tier mobile hardware, embedded Internet of Things (IoT) media players, and legacy smart televisions. Furthermore, streaming high-bitrate content over volatile networks with fluctuating bandwidth demands a leaner, more agile paradigm for cryptographic asset protection.

To address this technological tension between robust security and execution efficiency, this paper introduces a lightweight, codec-aware DRM architecture designed for secure and rapid video distribution. Instead of executing indiscriminate cryptographic operations across billions of continuous media bytes, the proposed system exploits the underlying syntactic structure of the H.264/Advanced Video Coding (AVC) bitstream. The architecture isolates and encrypts only high-priority structural anchoring components, specifically the Network Abstraction Layer (NAL) units corresponding to Instantaneous Decoder Refresh (IDR) frames along with critical sequence and picture parameter sets (SPS/PPS) \cite{winkler2024selective}. By deploying the Advanced Encryption Standard (AES) operating in Counter (CTR) mode exclusively on these critical references, the system completely breaks temporal and spatial decoding dependencies. The remaining bulk video data (P and B frame slices) is left unencrypted but remains entirely unwatchable and structurally unconstructible to an adversary.

The primary contributions of this work, which fulfill the empirical and design requirements are summarized as follows:
\begin{enumerate}
    \item The architectural layout of a modular, lightweight DRM system consisting of a secure offline packaging pipeline, a decoupled token-based license server, and an efficient client-side player.
    \item A codec-aware Selective Encryption (SE) mechanism utilizing AES-CTR that minimizes byte-mutation overhead while maintaining cryptographic resistance against visual reconstruction attacks.
    \item A comprehensive experimental profile assessing processing latency, CPU and memory utilization, and visual degradation metrics to empirically demonstrate the system's lightweight claim.
\end{enumerate}